\chapter{Extreme Weather}
\index{Extreme Weather}

\section*{Definition and Overview}
Extreme weather includes heatwaves, cold snaps, bushfires, floods, cyclones, and severe storms, all of which can have serious effects on health. Heatwaves are one of the deadliest natural hazards, and extremes of heat and cold, along with bushfire smoke, floods, and storms, are associated with injury, illness, and death, particularly in vulnerable people. This chapter explains the health risks of extreme weather, who is most at risk, and how to prepare for and respond to these events.

\section*{Epidemiology}
Extreme weather events are becoming more frequent and severe with climate change. Heatwaves cause more deaths in many countries than all other natural hazards combined, and heat-related illness affects people of all ages. Bushfire smoke causes widespread respiratory illness, and floods and storms cause injuries, infections, and mental health problems. Older people, young children, people with chronic illness, and those living alone are at greatest risk.

\section*{Aetiology and Pathophysiology}
Extreme weather harms health in several ways.
\begin{itemize}
    \item \textbf{Heat:} high temperatures overload the body's cooling system, causing heat exhaustion and life-threatening heatstroke; the heart and kidneys are strained
    \item \textbf{Cold:} low temperatures cause hypothermia, worsen heart and lung disease, and increase falls
    \item \textbf{Bushfire smoke:} fine particles irritate the lungs, triggering asthma and worsening heart and lung disease
    \item \textbf{Floods:} cause drowning, injury, infections from contaminated water, and displacement
    \item \textbf{Storms and cyclones:} cause injury from debris and structural damage, and disrupt power and services
    \item \textbf{Mental health:} extreme weather causes stress, anxiety, trauma, and post-traumatic stress
    \item \textbf{Vulnerability:} age, chronic illness, poverty, and social isolation increase the risks
\end{itemize}

\section*{Clinical Features}
The health effects vary with the event.
\begin{itemize}
    \item Heat-related illness: headache, dizziness, nausea, cramps, confusion, and collapse
    \item Hypothermia: shivering, confusion, drowsiness, and slurred speech
    \item Smoke effects: coughing, wheezing, breathlessness, and eye irritation
    \item Flood-related: injuries, wound infections, gastroenteritis from contaminated water, and hypothermia
    \item Storm-related: traumatic injuries and burns
    \item Psychological effects: anxiety, sleep disturbance, and distress
\end{itemize}

\section*{Differential Diagnosis}
Symptoms during extreme weather may have other causes.
\begin{itemize}
    \item Chest pain and breathlessness may indicate a heart attack rather than smoke exposure
    \item Confusion in heat may be heatstroke or another medical emergency
    \item Fatigue and headache may be dehydration or infection
    \item The context of the weather event guides assessment
\end{itemize}

\section*{When to Seek Medical Care}
Anyone with symptoms of heatstroke, hypothermia, or severe breathing problems should seek urgent care. People with chronic conditions should follow their action plans and seek help early.

\textbf{Contact your local emergency services for:}
\begin{itemize}
    \item Heatstroke: high body temperature, confusion, seizures, or collapse in the heat
    \item Hypothermia: confusion, severe shivering, or drowsiness in the cold
    \item Severe breathlessness or chest pain
    \item Drowning or near-drowning
    \item Serious injury during a storm or flood
\end{itemize}

\section*{Diagnosis and Investigations}
Assessment during extreme weather focuses on the emergency.
\begin{itemize}
    \item \textbf{Heat:} body temperature, hydration, and consciousness
    \item \textbf{Cold:} body temperature and level of consciousness
    \item \textbf{Smoke and flood:} breathing, oxygen, and signs of infection
    \item \textbf{Monitoring:} vulnerable people are checked on by family, friends, and services
\end{itemize}

\section*{Management}
Management involves immediate treatment and broader planning.
\begin{itemize}
    \item \textbf{Heat:} move to a cool place, remove clothing, drink fluids, and cool actively for heatstroke while awaiting help
    \item \textbf{Cold:} warm gradually, with medical care for hypothermia
    \item \textbf{Smoke:} stay indoors, use air filters, and follow asthma plans
    \item \textbf{Floods:} avoid floodwater, and seek medical care for wounds and illness
    \item \textbf{Preparedness:} emergency kits, plans, and staying informed
    \item \textbf{Community support:} checking on neighbours, especially the vulnerable
\end{itemize}

\section*{Complications}
\begin{itemize}
    \item Death from heatstroke and hypothermia
    \item Heart attacks and strokes during heatwaves
    \item Worsening of asthma and lung disease from smoke
    \item Infections and injuries from floods
    \item Displacement and loss of homes
    \item Post-traumatic stress and other mental health effects
\end{itemize}

\section*{Prognosis and Outcomes}
Most people recover fully from the health effects of extreme weather with prompt care, and heatstroke and hypothermia are treatable when recognised early. The wider impacts on communities can be long-lasting, which is why preparation, early warning, and community support matter. With planning, the risks of extreme weather can be substantially reduced.

\section*{Prevention and Self-Care}
\begin{itemize}
    \item Prepare an emergency kit and a family plan
    \item Stay informed about warnings and follow official advice
    \item During heat: stay cool, hydrated, and check on vulnerable people
    \item During smoke: stay indoors and follow asthma plans
    \item Never drive or walk through floodwater
    \item Protect against storms by securing loose items
    \item Seek support for mental health after extreme weather events
\end{itemize}

\section*{Sources and Further Reading}
The following reputable sources provide further information for healthcare students and patients seeking reliable, current advice about extreme weather and health.
\begin{itemize}
    \item Bureau of Meteorology. \textit{Weather warnings and information.} https://www.bom.gov.au/
    \item Department of Health and Aged Care. \textit{Heat and health advice.} https://www.health.gov.au/
    \item Healthdirect Australia. \textit{Extreme heat and health.} https://www.healthdirect.gov.au/heatwaves
    \item NSW Health. \textit{Health advice for extreme weather.} https://www.health.nsw.gov.au/
    \item Asthma Australia. \textit{Managing asthma in smoke and storms.} https://asthma.org.au/
    \item National Asthma Council Australia. \textit{Smoke and respiratory health.} https://www.nationalasthma.org.au/
\end{itemize}
